\documentclass[10pt,a4paper]{amsart}
\usepackage[margin=19mm]{geometry}
\usepackage{amsmath,amssymb,mathtools}
\usepackage[scr=rsfs,cal=euler]{mathalfa}
\usepackage{xfrac}
\usepackage[unicode,hidelinks,bookmarks=false]{hyperref}
\newcommand{\ie}{i.e.\ }
\newcommand{\cf}{cf.\ }
\newcommand{\resp}{respectively\ }
\newcommand{\Subsection}{\subsection}
\providecommand{\nobreakdash}{\nobreak\hbox{-}}
\setlength{\emergencystretch}{2em}
\multlinegap0pt
\usepackage{amssymb}
\usepackage{xcolor}
\usepackage{mathtools}
\newtheorem{theorem}{Theorem}
\newtheorem{lemma}{Lemma}
\newtheorem{claim}{Claim}
\newcommand{\R}{\mathbb{R}}
\newcommand{\eps}{\varepsilon}
\title{Sharp estimates for eigenvalues of localization operators with applications to area laws}
\author{Aleksei Kulikov, Martin Dam Larsen}
\date{Source: arXiv:2603.23832v1 (CC BY 4.0)}
\begin{document}
\maketitle

\noindent\emph{Source and licence.} Sharp estimates for eigenvalues of localization operators with applications to area laws. Version arXiv:2603.23832v1 (CC BY 4.0), distributed by arXiv under the Creative Commons Attribution 4.0 International licence, http://creativecommons.org/licenses/by/4.0/, as recorded in the arXiv licence metadata for this exact version and shown on the abstract page https://arxiv.org/abs/2603.23832v1. Full licence notice: \texttt{licenses/arxiv-2603.23832-CC-BY-4.0.txt}.\par\smallskip

\noindent\textit{Editorial scope.} Counted unit: the article's theorem bounded (source0.tex lines 549--555). The article's complete original proof of it is delivered with it (source0.tex lines 1492--1515, one \texttt{proof} environment of 24 lines, which the article places away from the statement). No mathematics is written, repaired or re-derived: the statement and the proof are the article's own lines, in the article's own notation, and no retained line is substituted or re-flowed.  Retained uncounted as the source-local context the proof needs: lines 296--309; lines 306--308; lines 665--667; lines 844--853; lines 1294--1298; lines 1378--1380; lines 1457--1463.  Nothing was omitted from any retained span, so the package carries no omission record.  No retained line contains a citation, so the wrapper carries no bibliography and no bibliography entry.  No retained line contains a figure environment, an image-inclusion command or drawing code, so no image is delivered and the package declares no asset dependency.  Editorial formatting only, in addition: an AMS \texttt{amsart} preamble that restates the article's own declarations and macros used by the retained lines, namely the environments \texttt{theorem}, \texttt{lemma}, \texttt{claim} on separate counters, together with the standard packages those lines need.  The retained environments print in this document as Theorem 1, Lemma 1, Claim 1 in order of appearance, whereas the article prints the same items with the numbering of its own full document.  The complete native source of the article as distributed in the arXiv e-print for this exact version is bundled unchanged as \texttt{sources/197112/source0.tex}.
\begin{theorem}\label{two cubes}
Consider $A = B = [0, 1]^d$, $d\geq 1$, and $c\geq 2$. There exists $\alpha_d \geq 4$ such that
\begin{equation}\label{Lambda pm, upper, two cubes}
\Lambda_\eps^{\pm}(c A, B) \lesssim c^{d-1} \log(\tfrac{1}{\eps}) \log\left(\frac{\alpha_d \, c}{\log\left(\tfrac{1}{\eps}\right)}\right)
\end{equation}
uniformly for all $\alpha_d^{-c}< \eps < 1/2$. Moreover, if $\eps < c^{-\alpha_d}$ then we also have
\begin{equation}\label{Lambda pm, lower, two cubes}
\Lambda_\eps^{\pm}(c A, B) \gtrsim c^{d-1} \log(\tfrac{1}{\eps}) \log\left(\frac{\alpha_d \, c}{\log\left(\tfrac{1}{\eps}\right)}\right)
\end{equation}
If $\eps \leq \alpha_d^{-c}$ then there are no eigenvalues larger than $1-\eps$ and 
\begin{equation}\label{Lambda -, equivalence, two cubes}
\Lambda_\eps^- \asymp \left( \frac{\log(\tfrac{1}{\eps})}{\log \left(\tfrac{\log\left(\frac{1}{\eps}\right)}{c}\right)}\right)^d.
\end{equation}
\end{theorem}

\begin{equation}
\Lambda_\eps^- \asymp \left( \frac{\log(\tfrac{1}{\eps})}{\log \left(\tfrac{\log\left(\frac{1}{\eps}\right)}{c}\right)}\right)^d.
\end{equation}

\begin{equation}\label{N 1/2}
N_{1/2}([0,c]^d,[0,1]^d) = c^d +O(c^{d-1}\log c).
\end{equation}

\begin{lemma}\label{enlargement}
Consider $X, Y, Z \subseteq \R^d$ and assume that $P_X Q_Y P_Z$ is compact. If $X\subset X'$ and $P_{X'}Q_Y P_Z$ is compact, then
\[
\sigma_k(P_X Q_Y P_Z) \leq \sigma_k(P_{X'} Q_Y P_Z).
\]
Similarly, if $Z\subseteq Z'$ and $P_X Q_Y P_{Z'}$ is compact, then 
\[
\sigma_k(P_X Q_Y P_Z) \leq \sigma_k(P_{X} Q_Y P_{Z'}).
\]
\end{lemma}

\begin{equation}\label{trace class condition bound}
\begin{aligned}\sum_{\alpha^{-c}>\lambda_n(cA,B) >0}|f(\lambda_n(cA,B))| &\lesssim \log(c)^d\sum_{k=k_0(c)}^\infty  M_0 f\left(2^{-2^k}\right) \frac{2^{(k+1)d}}{(k+1)^d}\\
&\lesssim \log(c)^d\int_0^{2^{-2^{k_0(c)-1}}}\frac{M_0 f(\eps)\log\left(\frac{1}{\eps}\right)^{d-1}}{\eps \left(\log\log\left(\frac{1}{\eps}\right)\right)^d} \, d\eps.
\end{aligned}
\end{equation}

\begin{equation}\label{cont at 0}
\max(M_0f(\eps), M_1f(\eps))\to 0 \text{ as } \eps \to 0^+.
\end{equation}

\begin{claim}\label{the trivial claim}
Let $A, B \subseteq \R^d$ be sets with finite measure. Then $S_{cA,B}^2$ is trace class and 
\[
\operatorname{Tr}S_{cA,B}^2 = c^d |A||B| + o(c^d) 
\]
as $c\to \infty$.
\end{claim}

\begin{theorem}\label{one term bounded}
Let $A, B \subseteq \R^d$ be bounded sets and assume that $f$ is bounded, trace class admissible for $L^2(\R^d)$ and that $f$ is continuous at $1$. Then 
\[
\operatorname{Tr} f(S_{cA, B}) = c^d|A||B| f(1) + o(c^{d}),
\]
as $c\to \infty$.
\end{theorem}

\begin{proof}[Proof of Theorem \ref{one term bounded}]
As in the previous proof we can without loss of generality assume that $f(1) = 0$. We also remark that since $f$ is trace class admissible for $L^2(\R^d)$, $f$ must satisfy $\lim_{\theta\to 0^+} f(\theta) =0$, by an argument similar to the proof of \eqref{cont at 0}, just with the trace class admissible for $L^2(\R^d)$ condition instead of the area law admissible condition.

Since $A, B$ are bounded, there exist boxes $V_A, V_B$ such that $A\subseteq V_A$, $B\subseteq V_B$. Fix once and for all a large number $D$ to be determined later (it will only depend on $V_A, V_B$). Let $0 < \eps < \frac{1}{2}$ be a small number. We have
\begin{align*}
{\rm Tr}f(S_{cA,B}) &= \sum_{\lambda_n(cA, B) > 1-\eps}f(\lambda_n(cA,B))  +\sum_{1-\eps\ge\lambda_n(cA, B) \geq \eps }f(\lambda_n(cA,B))\\
&\quad +\sum_{\substack{0 \le \lambda_n(cA,B)< \eps \\ n\le Dc^d}} f(\lambda_n(cA,B))+ \sum_{\substack{0 \le \lambda_n(cA,B)< \eps\\ n> Dc^d}} f(\lambda_n(cA,B)).
\end{align*}
For the first sum for any $\delta > 0$ we can choose $\frac{1}{2}>\eps > 0$ small enough so that it is at most $2\delta |A| |B| c^d$, as in the previous proof. 

For the second sum we simply use that $f$ is bounded and the elementary inequality $1\leq \frac{\theta(1-\theta)}{\eps(1-\eps)}$ for $\theta\in [\eps, 1-\eps]$. Thus, the second sum is at most $\frac{ {\rm Tr}[S_{cA,B}-S_{cA,B}^2]}{\eps(1-\eps)}\sup_{t\in [0,1]} |f(t)|$, which is $o(c^d)$ by Claim \ref{the trivial claim}. 

For the third sum, we use that $\lim_{\theta\to 0^+} f(\theta) =0$ to conclude that for any $\delta > 0$ we can choose $\eps > 0$ small enough so that it is at most $\delta Dc^d$.

Finally, we turn to the most challenging fourth sum. Firstly, we estimate $|f(\lambda_n(cA, B))|$ by $M_0f(\lambda_n(cA,B))$. Since $M_0f$ is non-decreasing, we can upper bound $\lambda_n(cA,B)$ by $\lambda_n(cV_A, V_B)$ using Lemma \ref{enlargement}. So, it remains to show that
$$\sum_{n > Dc^d} M_0f(\lambda_n(cV_A, V_B)) = o(c^d).$$
We claim that if $D$ is large enough and $n> D c^d$, then $\lambda_n(cV_A, V_B) \le \alpha_d^{-c}$, where $\alpha_d$ is taken from Theorem \ref{two cubes}. This is equivalent to saying that $N_{\alpha_d^{-c}}(cV_A, V_B)\le Dc^d$. We have
$$N_{\alpha_d^{-c}}(cV_A, V_B) = N_{1/2}(cV_A,V_B) + \Lambda_{\alpha_d^{-c}}(cV_A,V_B).$$
The first term is $O(c^d)$ by \eqref{N 1/2} and the second term is $O(c^d)$ by \eqref{Lambda -, equivalence, two cubes}. Thus, $N_{\alpha_d^{-c}}(cV_A, V_B)\le Dc^d$ holds for large enough $D$. Given the claim, the computation \eqref{trace class condition bound} therefore shows that
\[
\sum_{n > Dc^d} M_0f(\lambda_n(cV_A, V_B)) \leq \sum_{\lambda_n(cV_a,V_B) \leq \alpha_d^{-c}}  M_0f(\lambda_n(cV_A, V_B)) = O(\log(c)^d) = o(c^d),
\]
which finishes the proof.
\end{proof}

\end{document}
